\chapter{Tangent cones and stratification}\label{sec:tangent-cones-strata}
A blow-up records the geometry visible at smaller and smaller scales. Compactness provides limits, monotonicity forces them to be cones, and stratification measures their translational symmetries. Throughout this section $V$ is a stationary integral $n$-varifold in an open $U\subset\mathbb R^{n+k}$, where $n,k\geq1$.

\begin{definition}[Dilations of a varifold]\label{def:varifold-dilation}
For $a\in U$ and $r>0$, set $\eta_{a,r}(x)=(x-a)/r$. The dilated varifold $V_{a,r}=(\eta_{a,r})_\#V$ is defined on $\eta_{a,r}(U)$ by
\[
 \int\varphi(y,S)\,dV_{a,r}(y,S)
 =r^{-n}\int\varphi((x-a)/r,S)\,dV(x,S).
\]
The factor $r^{-n}$ is the tangential Jacobian of the dilation. This is not the unweighted pushforward of the measure on the Grassmann bundle. See \cite{Minter}, Section 1, p.~11.
\end{definition}

\begin{proposition}[Mass and stationarity under dilation]\label{prop:varifold-scaling}
For every ball with $B_{rR}(a)\subset U$,
\[
 \|V_{a,r}\|(B_R(0))=r^{-n}\|V\|(B_{rR}(a)).
\]
Moreover $V_{a,r}$ is stationary and integral, and $\Theta^n(\|V_{a,r}\|,0)=\Theta^n(\|V\|,a)$.
\end{proposition}
\begin{proof}
The mass identity follows by testing the definition with approximations to the indicator of $B_R$. If $Y$ is a compactly supported vector field in the dilated domain, put $X(x)=Y((x-a)/r)$. Then
$\delta V_{a,r}(Y)=r^{1-n}\delta V(X)=0$.
Rectifiable sets and integer multiplicities are preserved by the dilation. Divide the mass identity by $\omega_n R^n$ and send $R$ to zero to obtain the density identity.
\end{proof}

\begin{proposition}[Upper semicontinuity of stationary density]\label{prop:density-upper-semicontinuity}
If stationary integral varifolds $V_j$ converge to $V$ in $U$ and $a_j\to a\in U$, then
\[
 \limsup_j\Theta^n(\|V_j\|,a_j)\leq\Theta^n(\|V\|,a).
\]
In particular the density of a fixed stationary integral varifold is upper semicontinuous in the center. Compare \cite{Minter}, Lemma 0.6.
\end{proposition}
\begin{proof}
Fix $s>0$ with $\overline{B_s(a)}\subset U$ and $\|V\|(\partial B_s(a))=0$. For $0<\varepsilon<s$ and all large $j$, $B_{s-\varepsilon}(a_j)\subset B_s(a)$. Monotonicity bounds the density at $a_j$ by $\|V_j\|(B_s(a))/(\omega_n(s-\varepsilon)^n)$. At the chosen continuity radius the numerator converges to $\|V\|(B_s(a))$. Take the upper limit, then let $\varepsilon\downarrow0$, and finally let $s\downarrow0$ through continuity radii. Only countably many radii can fail the boundary-mass condition.
\end{proof}

\begin{corollary}[Density at every support point]\label{cor:integral-support-density}
For every $a\in\operatorname{spt}\|V\|$, $\Theta^n(\|V\|,a)\geq1$.
\end{corollary}
\begin{proof}
Integrality gives positive integer density at $\|V\|$-almost every point. Every ball about a support point has positive weight, so choose such points tending to $a$. Proposition \ref{prop:density-upper-semicontinuity}, applied to the constant sequence $V$, gives the lower bound. This conclusion uses both integrality and stationarity.
\end{proof}

\begin{definition}[Regular and singular points]\label{def:varifold-regular-singular}
A support point $a$ is regular if in some neighborhood $V=q|\Sigma|$, where $\Sigma$ is a smooth embedded $n$-submanifold without boundary there and $q$ is a positive integer. The remaining support points form $\operatorname{sing}V$; the regular set is $\operatorname{reg}V$. Locally constant multiplicity is included in regularity, so a multiple plane is regular. Transverse crossings are singular for this definition. For stationary integral varifolds this agrees with the smooth-support definition using Theorem \ref{thm:varifold-smooth-constancy}; see \cite{Minter}, Lemma 0.7 and Definitions 0.8--0.9.
\end{definition}

\begin{theorem}[Existence and structure of tangent cones]\label{thm:tangent-cones}
Let $a\in\operatorname{spt}\|V\|$. For every sequence $r_j\downarrow0$, a subsequence of $V_{a,r_j}$ converges locally as varifolds on $\mathbb R^{n+k}$ to a nonzero stationary integral $n$-varifold $C$. The limit is a cone:
\[
 (\eta_{0,s})_\#C=C\quad(s>0),\qquad
 \|C\|(B_R(0))=\omega_n R^n\Theta^n(\|V\|,a)\quad(R>0).
\]
Such limits are called tangent cones; their set is denoted $\operatorname{Tan}(V,a)$. No uniqueness is asserted. Sources: \cite{Minter}, Section 1, Lemmas 1.1--1.2 and Definition 1.3. The compactness input remains external; the conical-structure deduction follows below.
\end{theorem}
\begin{proof}
Choose $r_0>0$ with $\overline{B_{r_0}(a)}\subset U$. For fixed $R$ and all sufficiently large $j$, scaling and stationary monotonicity give
\[
 \|V_{a,r_j}\|(B_R)=r_j^{-n}\|V\|(B_{r_jR}(a))
 \leq R^n r_0^{-n}\|V\|(B_{r_0}(a)).
\]
The dilates are integral and stationary by Proposition \ref{prop:varifold-scaling}; hence their first variations vanish on every fixed compact set for large $j$. Allard compactness, Theorem \ref{thm:allard-compactness}, on an exhaustion and a diagonal subsequence give a stationary integral limit $C$ on all of $\mathbb R^{n+k}$. Stationarity also follows directly from Lemma \ref{lem:stationarity-closed}.

At any radius $R$ with $\|C\|(\partial B_R)=0$, weak convergence and scaling imply
\[
 \|C\|(B_R)=\lim_j r_j^{-n}\|V\|(B_{r_jR}(a))
            =\omega_nR^n\Theta^n(\|V\|,a).
\]
Only countably many radii fail this continuity condition. Approximating any radius from below by continuity radii and using continuity from below of open-ball masses proves the identity for every $R>0$. Theorem \ref{thm:constant-ratio-cone} now proves dilation invariance, including the tangent-plane distribution. Corollary \ref{cor:integral-support-density} gives $\Theta^n(\|V\|,a)\geq1$, so the limit is nonzero. The construction uses subsequences and supplies no uniqueness statement.
\end{proof}

\begin{definition}[The spine of a stationary cone]\label{def:cone-spine}
For a nonzero stationary integral cone $C$ in $\mathbb R^{n+k}$, its spine is the set of translation symmetries
\[
 S(C)=\{y\in\mathbb R^{n+k}:(\eta_{y,1})_\#C=C\}.
\]
Translation acts on the entire varifold, including multiplicity. See \cite{Minter}, Definition 1.4.
\end{definition}

\begin{theorem}[Spine structure and splitting]\label{thm:cone-spine-splitting}
The spine of a nonzero stationary integral $n$-cone is a linear subspace of dimension $d\leq n$, and
\[
 S(C)=\{y:\Theta^n(\|C\|,y)=\Theta^n(\|C\|,0)\}.
\]
The vertex has maximal density. After an orthogonal change of coordinates, $C=C_0\times\mathbb R^d$, where $C_0$ is a stationary integral $(n-d)$-cone; for $d=n$, $C_0$ is a positive integer multiple of the origin. Source: \cite{Minter}, Lemma 1.5.
\end{theorem}
\begin{proof}
Proposition \ref{prop:spine-maximal-density} identifies the maximal-density set with a linear spine of dimension at most $n$. If $d=0$, take $C_0=C$. If $1\leq d<n$, Proposition \ref{prop:cone-product-splitting} gives the stationary integral transverse cone, using rectifiable coarea as its recorded external input. If $d=n$, Corollary \ref{cor:full-spine-plane} gives a positive integer multiple of the spine, with transverse factor $q\delta_0$. These statements identify the weight, tangent planes, multiplicity, and first variation of each factor.
\end{proof}

\begin{definition}[Singular strata]\label{def:singular-strata}
For $0\leq j\leq n$, define
\[
 \mathcal S^j(V)=\{a\in\operatorname{sing}V:
 \dim S(C)\leq j\text{ for every }C\in\operatorname{Tan}(V,a)\}.
\]
Thus $\mathcal S^0\subset\cdots\subset\mathcal S^n=\operatorname{sing}V$. The quantifier is over every tangent cone, without assuming tangent-cone uniqueness. A point with a planar tangent cone of higher multiplicity can lie in the top stratum. See \cite{Minter}, Definition 1.6.
\end{definition}

\begin{theorem}[Dimension bound for the strata]\label{thm:strata-dimension}
For every stationary integral $n$-varifold and $0\leq j\leq n$,
\[
 \dim_{\mathcal H}\mathcal S^j(V)\leq j.
\]
This is the qualitative stratification theorem stated in \cite{Minter}, Theorem 1.7, p.~17. Its covering proof is given below. The estimate alone does not remove the top stratum or give the codimension-seven bound available under stronger minimizing or stability hypotheses.
\end{theorem}
\begin{proof}
Fix $\beta>0$ and choose $\delta>0$ small enough for Proposition \ref{prop:flat-approximation-dimension} in ambient dimension $n+k$ and plane dimension $j$. Proposition \ref{prop:stratum-flat-pieces}, proved by tangent-cone compactness and density upper semicontinuity, decomposes $\mathcal S^j(V)$ into countably many sets with one-sided $\delta$ approximation by $j$-planes at all sufficiently small scales. Each piece has zero $\mathcal H^{j+\beta}$ outer measure by the covering proposition. Countable subadditivity gives the same conclusion for the stratum. Since $\beta>0$ is arbitrary, its Hausdorff dimension is at most $j$. This proves a dimension bound; it does not prove finite $\mathcal H^j$ measure, rectifiability, or uniqueness of tangent cones.
\end{proof}


\begin{example}[Crossing planes remain stationary]\label{ex:crossing-planes}
Let $P_1,P_2$ be distinct $n$-dimensional hyperplanes in $\mathbb R^{n+1}$ through the origin. The varifold $|P_1|+|P_2|$ is stationary by linearity and Proposition \ref{prop:plane-stationary}. Its singular set and spine are $P_1\cap P_2$, of dimension $n-1$. The density is two on the intersection and one elsewhere on the support. Thus stationarity and integrality allow a large singular set. Compare \cite{Minter}, Sections 0--1.
\end{example}

\chapter{Excess and Allard regularity}\label{sec:allard-regularity}
Being close to a plane has several meanings. Height excess measures distance from the plane, tilt excess measures tangent directions, and a mass-ratio bound controls multiplicity. These quantities cannot be substituted for one another without a theorem.

\begin{definition}[Height and tilt excess]\label{def:varifold-excess}
Let $V$ be an $n$-varifold in $B_R(a)\subset\mathbb R^{n+k}$ and $P$ a linear $n$-plane. With $\pi_S$ denoting orthogonal projection onto $S$, set
\[
 E_h(V,a,R,P)^2=R^{-n-2}\int_{B_R(a)}\operatorname{dist}(x,a+P)^2\,d\|V\|(x),
\]
\[
 E_t(V,a,R,P)^2=R^{-n}\int_{G_n(B_R(a))}|\pi_S-\pi_P|_{\mathrm{HS}}^2\,dV(x,S).
\]
Here $|\cdot|_{\mathrm{HS}}$ is the Hilbert--Schmidt norm. These are dimensionless quantities; our tilt normalization is fixed by this formula and need not equal another author's scalar-normal convention. Compare \cite{Sim18}, Chapter 5, and \cite{Minter}, Sections 2--3.
\end{definition}

\begin{proposition}[Excess is invariant under simultaneous rescaling]\label{prop:excess-scaling}
For $V_{a,R}$ as in Definition \ref{def:varifold-dilation},
\[
 E_h(V,a,R,P)=E_h(V_{a,R},0,1,P),\qquad
 E_t(V,a,R,P)=E_t(V_{a,R},0,1,P).
\]
If the first variation is entirely represented by a mean-curvature field $H\in L^p(\|V\|)$, its rescaled field is $H_R(y)=R H(a+Ry)$, and
\[
 \left(\int_{B_1}|H_R|^p\,d\|V_{a,R}\|\right)^{1/p}
 =R^{1-n/p}\left(\int_{B_R(a)}|H|^p\,d\|V\|\right)^{1/p}.
\]
\end{proposition}
\begin{proof}
The mass factor is $R^{-n}$, distances scale by $R^{-1}$, and planes are unchanged. Substitute into the excess definitions. For the curvature formula, use the first-variation scaling from Proposition \ref{prop:varifold-scaling} and the identity $\delta V(X)=-\int H\cdot X\,d\|V\|$. The last equality is then a change of variables.
\end{proof}

\begin{theorem}[Allard regularity with supercritical mean curvature]\label{thm:allard-lp}
Fix $n,k\geq1$ and $n<p<\infty$. There are $\varepsilon>0$ and $\gamma\in(0,1)$, depending only on $n,k,p$, with the following property. Let $V$ be an integral $n$-varifold in an open neighborhood of $\overline{B_R(a)}\subset\mathbb R^{n+k}$, with $a\in\operatorname{spt}\|V\|$, and suppose its first variation in $B_R(a)$ has no singular part:
\[
 \delta V(X)=-\int H\cdot X\,d\|V\|\quad
 (X\in C_c^1(B_R(a);\mathbb R^{n+k})).
\]
Assume $H\in L^p(\|V\|\llcorner B_R(a))$ and
\[
 \frac{\|V\|(B_R(a))}{\omega_n R^n}\leq1+\varepsilon,
 \qquad R^{1-n/p}\left(\int_{B_R(a)}|H|^p\,d\|V\|\right)^{1/p}\leq\varepsilon.
\]
Then in $B_{\gamma R}(a)$ the support of $V$ is a single $C^{1,1-n/p}$ graph over an $n$-plane, and $V$ has multiplicity one there. This is the integral-varifold specialization of \cite{Sim18}, Chapter 5, hypotheses (5.1) and Theorem 5.2, pp.~145--146, following \cite{Allard72}. The source's density-lower-bound hypothesis is supplied here by integrality almost everywhere. The proof, including affine approximation and excess decay, remains external. General bounded first variation with a singular boundary term does not satisfy the stated hypothesis.
\end{theorem}

\begin{corollary}[Density one forces a regular point]\label{cor:unit-density-regularity}
If $V$ is stationary and integral and $\Theta^n(\|V\|,a)=1$, then $a$ is regular with multiplicity one. In particular one multiplicity-one planar tangent cone is enough to imply regularity and uniqueness of the tangent cone at $a$.
\end{corollary}
\begin{proof}
At a sufficiently small radius the mass ratio is at most $1+\varepsilon$, while the mean-curvature term is zero. Theorem \ref{thm:allard-lp} gives a $C^{1,\alpha}$ single graph. The stationary graph equation and elliptic bootstrapping make this graph smooth; this stationary consequence is also recorded in \cite{Minter}, Corollary 3.2. A multiplicity-one planar tangent cone has density one by Theorem \ref{thm:tangent-cones}. At a smooth point every blow-up is the ordinary tangent plane with its multiplicity.
\end{proof}

\begin{theorem}[A density gap above one]\label{thm:stationary-density-gap}
There is $\delta=\delta(n,k)>0$ such that every support point $a$ of a stationary integral $n$-varifold in $\mathbb R^{n+k}$ with $\Theta^n(\|V\|,a)<1+\delta$ is regular with multiplicity one. Equivalently, every singular point has density at least $1+\delta$. This is \cite{Minter}, Corollary 3.2, pp.~44--45. Its Allard regularity input remains external.
\end{theorem}
\begin{proof}
Let $\varepsilon=\varepsilon(n,k,2n)$ be the mass tolerance in Theorem \ref{thm:allard-lp}, and take $\delta=\varepsilon/2$. At a point with density less than $1+\delta$, the mass ratio is less than $1+\varepsilon$ at all sufficiently small radii. Stationarity gives $H=0$, so the curvature hypothesis holds with $p=2n>n$. Allard regularity gives a multiplicity-one $C^{1,1/2}$ graph near the point. The stationary graph equation and elliptic bootstrapping make it smooth, as in Corollary \ref{cor:unit-density-regularity}. Thus the gap is a local consequence of the already stated Allard theorem.
\end{proof}


\begin{example}[Zero excess does not specify multiplicity]\label{ex:multiple-plane-excess}
For $V=q|P|$ with integer $q\geq2$, both excesses relative to $P$ vanish at every scale centered on $P$, but the mass ratio and density are $q$. Thus zero height and tilt excess do not imply a single sheet with multiplicity one. This varifold is nevertheless regular under Definition \ref{def:varifold-regular-singular}; high multiplicity and singular support are different issues.
\end{example}
